\everymath{\displaystyle}

%%%%%%%%%%%%%%%%%%%%%%%%%%%%%%%%%%%%%%%%
%           main body
%%%%%%%%%%%%%%%%%%%%%%%%%%%%%%%%%%%%%%%%

%-----------------------------------
%	TITLE PAGE
%-----------------------------------

\title[第十章\\
函数项级数]{\texorpdfstring{\LARGE \bfseries 第十章\quad 函数项级数}{第十章 函数项级数}} % The short title appears at the bottom of every slide, the full title is only on the title page
\author[\smaller \S 10.3\\
幂级数] % Your institution as it will appear on the bottom of every slide, maybe shorthand to save space
{\Large \bfseries \S 10.3 幂级数}
% \author{Singular Value Decomposition} % Your name
% \date{\today} % Date, can be changed to a custom date
\date{}

%------------------------------------------------

\begin{document}

%------------------------------------------------

\begin{frame}

\titlepage % Print the title page as the first slide

\end{frame}

%-----------------------------------
%	PRESENTATION SLIDES
%-----------------------------------

%------------------------------------------------

\begin{frame}
\frametitle{回顾}

我们之前考虑过如下的例子:

设 $\displaystyle \sum_{n=0}^{\infty} a_n$ 收敛, 考虑函数项级数 $\displaystyle \sum_{n=0}^{\infty} a_n x^n:$

\begin{itemize}
\item 问题一: 它的收敛域 $D$ 是什么? \pause ~ $\rightsquigarrow$ $(-1, 1] \subset D.$
\pause
\item 问题二: 在收敛域上是否 (内闭) 一致收敛?
\end{itemize}

\pause
\vspace{1em}

这一节我们主要考察具有上述特殊形式 (通项为单项式) 的函数项级数, 并回答以上两个问题.

\pause

\begin{Def}[幂级数]
称形如 $\displaystyle \sum_{n=0}^{\infty} a_n (x - x_0)^n$ 的函数项级数为幂级数.
\end{Def}

\alert{注意}, 我们通常考虑 $x_0 = 0$ 的幂级数, 因为 $x \mapsto x - x_0$ 即可简单地将相关结论推广至 $x_0 \neq 0$ 的情形.

\end{frame}

%------------------------------------------------

\section{幂级数的收敛半径}

%------------------------------------------------

\begin{frame}
\frametitle{Cauchy-Hadamard 定理}

幂级数 $\displaystyle \sum_{n=0}^{\infty} a_n x^n$ 对于固定的 $x,$ 由正项级数敛散性的 Cauchy 判别法, 有
\begin{equation*}
\varlimsup_{n\to\infty} \nroot[n]{|a_n x^n|} = |x| \cdot \varlimsup_{n\to\infty} \nroot[n]{|a_n|} ~ \begin{cases}
< 1 ~ \text{时 (绝对) 收敛}, \\ > 1 ~ \text{时发散}.
\end{cases}
\end{equation*}
于是, $\displaystyle |x| < 1 \left/ \varlimsup_{n\to\infty} \nroot[n]{|a_n|} \right.$ 时, 幂级数 $\displaystyle \sum_{n=0}^{\infty} a_n x^n$ (绝对) 收敛. 即有

\pause

\begin{thm}[Cauchy-Hadamard 定理]
令 $\displaystyle A = \varlimsup_{n\to\infty} \nroot[n]{|a_n|},$ $R = \begin{cases}
+\infty, & A = 0, \\ 1/A, & A \in (0, +\infty), \\ 0, & A = +\infty,
\end{cases}$ 则幂级数 $\displaystyle \sum_{n=0}^{\infty} a_n x^n$ 当 $|x| < R ~ (R > 0)$ 时绝对收敛; 当 $|x| > R$ 时发散. $R$ 称为幂级数的\alert{收敛半径}.
\end{thm}

\end{frame}

%------------------------------------------------

\begin{frame}
\frametitle{例题}

\begin{eg}
分别求幂级数 $\displaystyle \sum_{n=1}^{\infty} \dfrac{x^n}{n}$ 与 $\displaystyle \sum_{n=1}^{\infty} \dfrac{(-1)^{n+1}}{n}x^n$ 的收敛半径以及收敛域.
\end{eg}

\pause

\startproof[bg=yellow, fg=red](解)
(1) $\displaystyle \sum_{n=1}^{\infty} \dfrac{x^n}{n}$ 的收敛半径为
\[
1 \left/ \varlimsup_{n\to\infty} \sqrt[\leftroot{-3}\uproot{7}n]{\frac{1}{n}} \right. = 1
\]
\pause
分别代入 $x = -1, 1$ 可得收敛域为 $[-1, 1).$

\pause

(2) $\displaystyle \sum_{n=1}^{\infty} \dfrac{(-1)^{n+1}}{n}x^n$ 的收敛半径为
\[
1 \left/ \varlimsup_{n\to\infty} \sqrt[\leftroot{-3}\uproot{7}n]{\frac{1}{n}} \right. = 1
\]
\pause
收敛域为 $(-1, 1].$

\end{frame}

%------------------------------------------------

\begin{frame}
\frametitle{例题}

\begin{eg}
求幂级数 $\displaystyle \sum_{n=1}^{\infty} \dfrac{x^n}{n!}$ 的收敛半径以及收敛域.
\end{eg}

\pause

\startproof[bg=yellow, fg=red](解)
$\displaystyle \sum_{n=1}^{\infty} \dfrac{x^n}{n!}$ 的收敛半径为
\[
1 \left/ \varlimsup_{n\to\infty} \sqrt[\leftroot{-3}\uproot{7}n]{\frac{1}{n!}} \right. = +\infty,
\]
\pause
所以收敛域等于 $\mathbb{R}.$

\end{frame}

%------------------------------------------------

\begin{frame}
\frametitle{例题}

\begin{eg}
考察幂级数 $\displaystyle \sum_{n=1}^{\infty} \dfrac{(2 + (-1)^n)^n}{n} \left( x - \frac{1}{2} \right)^n$ 的收敛情况.
\end{eg}

\pause
\vspace{0.5em}

\startproof[bg=yellow, fg=red](解)
记 $a_n = \dfrac{(2 + (-1)^n)^n}{n},$ 则 $\displaystyle \varlimsup_{n\to\infty} \sqrt[\leftroot{-3}\uproot{7}n]{|a_n|} = 3,$ 所以收敛半径 $R = 1/3.$

\pause
\vspace{0.5em}

$1^{\circ}$ 当 $x - 1/2 = 1/3$ 时, 得到一个正项级数 $\displaystyle \sum_{n=1}^{\infty} a_n \left( \frac{1}{3} \right)^n = \sum_{n=1}^{\infty} \frac{(2 + (-1)^n)^n}{n \cdot 3^n} =: \sum_{n=1}^{\infty} b_n.$
\[
b_n = \begin{cases}
\frac{1}{2m}, & n = 2m, ~ \rightsquigarrow \text{发散}, \\ \frac{1}{(2m - 1) \cdot 3^{2m - 1}}, & n = 2m - 1, ~ \rightsquigarrow \text{收敛}.
\end{cases}
~~ \implies \text{发散}.
\]

\end{frame}

%------------------------------------------------

\begin{frame}
\frametitle{例题}

$2^{\circ}$ 当 $x - 1/2 = -1/3$ 时, 得到 $\displaystyle \sum_{n=1}^{\infty} a_n \left( -\frac{1}{3} \right)^n = \sum_{n=1}^{\infty} \frac{(-2 + (-1)^{n+1})^n}{n \cdot 3^n} =: \sum_{n=1}^{\infty} c_n.$
\[
c_n = \begin{cases}
\frac{1}{2m}, & n = 2m, ~ \rightsquigarrow \text{发散}, \\ \frac{-1}{(2m - 1) \cdot 3^{2m - 1}}, & n = 2m - 1, ~ \rightsquigarrow \text{收敛}.
\end{cases}
~~ \implies \text{发散}.
\]

\pause

综上, 幂级数 $\displaystyle \sum_{n=1}^{\infty} \dfrac{(2 + (-1)^n)^n}{n} \left( x - \frac{1}{2} \right)^n$ 的收敛域为 $\left( \frac{1}{6}, \frac{5}{6} \right).$

\end{frame}

%------------------------------------------------

\begin{frame}
\frametitle{d'Alembert 判别法}

正项级数的敛散性判别法除了 Cauchy 判别法之外, 还有 d'Alembert 判别法. 对应地, 幂级数的收敛半径也可以通过 d'Alembert 判别法来计算:

\begin{thm}[d'Alembert 判别法]
若对幂级数 $\displaystyle \sum_{n=0}^{\infty} a_n x^n$ 成立 $\displaystyle \lim_{n\to\infty} \left| \frac{a_{n+1}}{a_n} \right| = A,$ 则收敛半径 $R = 1/A.$
\end{thm}

\pause

\startproof[bg=yellow, fg=red]

\end{frame}

%------------------------------------------------

\section{幂级数的分析性质}

%------------------------------------------------

\begin{frame}

Abel 曾系统地研究过幂级数, 对于本节一开始提的两个问题
\begin{itemize}
\item 幂级数的收敛域
\item 幂级数在收敛域内的分析性质
\end{itemize}
建立了一系列的定理.

\end{frame}

%------------------------------------------------

\begin{frame}
\frametitle{Abel 第一定理}

\begin{thm}[Abel 第一定理]

\end{thm}

\end{frame}

%------------------------------------------------

\begin{frame}
\frametitle{Abel 第二定理}

\begin{thm}[Abel 第二定理]

\end{thm}

\end{frame}

%------------------------------------------------

\begin{frame}
\frametitle{级数求和的 Abel 方法}

之前, 我们已经证明了: 若数项级数 $\displaystyle \sum_{n=1}^\infty a_n$ 收敛, 则幂级数 $\displaystyle \sum_{n=1}^\infty a_n x^n$ 的收敛半径 $R \geqslant 1,$ 且收敛域包含 $(-1, 1].$ 于是由 Abel 第二定理, $\displaystyle \sum_{n=1}^\infty a_n x^n$ 在 $[0, 1]$ 上一致收敛, 从而连续, 于是成立
\[
\sum_{n=1}^\infty a_n = \lim_{x \to 1^-} \sum_{n=1}^{\infty} a_n x^n
\]

\pause

但是, 有时当上式左边发散, 即级数 $\displaystyle \sum_{n=1}^\infty a_n$ 发散时, 右边 (左极限) 却是有意义的.
\pause
例如 $a_n = (-1)^n,$ 左边是一个发散的级数, 但
\begin{equation*}
\text{右边} = \lim_{x \to 1^-} \sum_{n=1}^{\infty} (-1)^n x^n = \lim_{x \to 1^-} \dfrac{1}{1 + x} = \dfrac{1}{2}.
\end{equation*}

\end{frame}

%------------------------------------------------

\begin{frame}
\frametitle{级数求和的 Abel 方法}

\begin{Def}[Abel 可和级数]
若极限 $\displaystyle \lim_{x \to 1^-} \sum_{n=1}^{\infty} a_n x^n = A$ 存在, 为有限实数, 则称数项级数 $\displaystyle \sum_{n=1}^\infty a_n$ 是 Abel 意义下可和的, 并称它的 Abel 意义下的和为 $A.$
\end{Def}

\pause
\vspace{1em}

\begin{prop}
若数项级数 $\displaystyle \sum_{n=1}^\infty a_n$ 在 Cesàro $(c, 1)$ 意义下的和为 $A \in \mathbb{R},$ 那么它在 Abel 意义下也是可和的, 且和也是 $A.$
\end{prop}

\pause

简单来说, 就是
\[
\{ \text{Cesàro $(c, 1)$ 可和级数} \} \subset \{ \text{Abel 可和级数} \}
\]

\end{frame}

%------------------------------------------------

\begin{frame}
\frametitle{$\{ \text{Cesàro $(c, 1)$ 可和级数} \} \subset \{ \text{Abel 可和级数} \}$ ~ 证明}

由于 $\displaystyle \sigma_n = \dfrac{1}{n} \sum_{k=1}^{n} S_k$ 极限为 $A,$ 所以
\[
\varlimsup_{n\to\infty} \nroot[n]{|n \sigma_n|} = \sqrt[\leftroot{-1}\uproot{18}n]{\sum_{k=1}^{n} S_k} = 1,
\]
即幂级数 $\displaystyle \sum_{n=1}^{\infty} n \sigma_n x^n$ 收敛半径为 $1.$
\pause
对于任意的 $|x| < 1,$ 上述幂级数绝对收敛, 因此有
\begin{align*}
\sum_{n=1}^{\infty} n \sigma_n x^n & = \sum_{n=1}^{\infty} \left( \sum_{k=1}^{n} S_k \right) x^n = \sum_{n=1}^{\infty} \left( S_1 + S_2 + \cdots + S_n \right) x^n \\
& = S_1 x (1 + x + x^2 + \cdots) + S_2 x^2 (1 + x + x^2 + \cdots) + \cdots \\
& = \sum_{n=1}^{\infty} S_n x^n \dfrac{1}{1-x}
\end{align*}

\end{frame}

%------------------------------------------------

\begin{frame}
\frametitle{$\{ \text{Cesàro $(c, 1)$ 可和级数} \} \subset \{ \text{Abel 可和级数} \}$ ~ 证明}

上式也表明了 $\displaystyle \sum_{n=1}^{\infty} S_n x^n$ 收敛半径 $\geqslant 1,$ 故在 $|x| < 1$ 范围内绝对收敛. 于是类似地有
\begin{align*}
\sum_{n=1}^{\infty} S_n x^n & = \sum_{n=1}^{\infty} \left( a_1 + a_2 + \cdots + a_n \right) x^n = \sum_{n=1}^{\infty} a_n x^n (1 + x + x^2 + \cdots) \\
& = \sum_{n=1}^{\infty} a_n x^n \dfrac{1}{1-x}.
\end{align*}
那么 $\displaystyle \sum_{n=1}^{\infty} n \sigma_n x^n = \dfrac{1}{(1-x)^2} \sum_{n=1}^{\infty} a_n x^n$ 对 $|x| < 1$ 恒成立.

\pause

接下来, 我们要证明
\[
\lim_{x\to 1-} (1-x)^2 \sum_{n=1}^{\infty} n \sigma_n x^n = \lim_{x\to 1-} \sum_{n=1}^{\infty} a_n x^n = A.
\]

\end{frame}

%------------------------------------------------

\begin{frame}
\frametitle{$\{ \text{Cesàro $(c, 1)$ 可和级数} \} \subset \{ \text{Abel 可和级数} \}$ ~ 证明}

对定义在 $|x| < 1$ 上的幂级数展开 $\displaystyle (1-x)^{-1} = \sum_{n=0}^\infty x^n$ 用应用逐项求导定理, 有$\displaystyle (1-x)^{-2} = \sum_{n=0}^\infty (n+1) x^n,$ 那么
\begin{align*}
& (1-x)^2 \sum_{n=1}^{\infty} n \sigma_n x^n - A \\
= & (1-x)^2 \sum_{n=1}^{\infty} n \sigma_n x^n - (1-x)^2 \sum_{n=0}^\infty (n+1) x^n A \\
= & (1-x)^2 \sum_{n=1}^{\infty} (n \sigma_n - (n+1)A) x^n - (1-x)^2 A \\
= & (1-x)^2 \sum_{n=1}^{\infty} n (\sigma_n - A) x^n - (1-x)^2 \sum_{n=1}^{\infty} A x^n - (1-x)^2 A
\end{align*}
上式后两项关于 $x \to 1-$ 的极限很容易看出都是 $0.$

\end{frame}

%------------------------------------------------

\begin{frame}
\frametitle{$\{ \text{Cesàro $(c, 1)$ 可和级数} \} \subset \{ \text{Abel 可和级数} \}$ ~ 证明}

我们最终约化到了证明第一项 $\displaystyle (1-x)^2 \sum_{n=1}^{\infty} n (\sigma_n - A) x^n$ 关于 $x \to 1-$ 的极限也是 $0.$

由于 $\sigma_n \to A,$ 所以 $\varepsilon > 0, ~ \exists N,$ 使得 $\forall~n > N,$ 有 $|\sigma_n - A| < \varepsilon / 2.$ 记 $\displaystyle M = \max_{1 \leqslant n \leqslant N} n |\sigma_n - A|,$ 那么有
{\smaller[1]
\begin{align*}
\left\lvert (1-x)^2 \sum_{n=1}^{\infty} n (\sigma_n - A) x^n \right\rvert & \leqslant \sum_{n=1}^{\infty} n |\sigma_n - A| (1-x)^2x^n \\
& = \sum_{n=1}^{N} n |\sigma_n - A| (1-x)^2x^n + \sum_{n=N+1}^{\infty} n |\sigma_n - A|(1-x)^2x^n \\
& \leqslant M \sum_{n=1}^{N} (1-x)^2x^n + \dfrac{\varepsilon}{2} \sum_{n=N+1}^{\infty} n (1-x)^2x^n \\
& \leqslant M x(1-x)(1-x^N) + \dfrac{\varepsilon}{2} \sum_{n=0}^{\infty} (n + 1) (1-x)^2x^n \\
& \leqslant M(1 - x) + \dfrac{\varepsilon}{2}.
\end{align*}
}

\end{frame}

%------------------------------------------------

\begin{frame}
\frametitle{$\{ \text{Cesàro $(c, 1)$ 可和级数} \} \subset \{ \text{Abel 可和级数} \}$ ~ 证明}

对取定的 $\varepsilon,$ 任取 $1 - \dfrac{\varepsilon}{2M} < x < 1,$ 即有
\[
\left\lvert (1-x)^2 \sum_{n=1}^{\infty} n (\sigma_n - A) x^n \right\rvert \leqslant M(1 - x) + \dfrac{\varepsilon}{2} < \varepsilon,
\]
这样, 我们就证明了
\[
\lim_{x\to 1-} \sum_{n=1}^{\infty} a_n x^n = \lim_{x\to 1-} (1-x)^2 \sum_{n=1}^{\infty} n \sigma_n x^n = A.
\]

\end{frame}

%------------------------------------------------

\begin{frame}
\frametitle{和函数连续性}

\end{frame}

%------------------------------------------------

\begin{frame}
\frametitle{逐项可积性}

\end{frame}

%------------------------------------------------

\begin{frame}
\frametitle{逐项可导性}

\end{frame}

%------------------------------------------------

\begin{frame}
\frametitle{关于数项级数 Cauchy 乘积的扩展结论}

之前讲数项级数的乘法时, 我们证明了若 $\displaystyle \sum\limits_{n=1}^{\infty} a_n, \sum\limits_{n=1}^{\infty} b_n$ 中至少一个绝对收敛时, 它们的 Cauchy 乘积 $\displaystyle \sum\limits_{n=2}^{\infty} {\color{magenta}\sum\limits_{k=1}^{n-1} a_k b_{n-k}} = \sum\limits_{n=2}^{\infty} {\color{magenta}c_n}$ 收敛, \pause
并且看到了 $\displaystyle \sum\limits_{n=1}^{\infty} a_n, \sum\limits_{n=1}^{\infty} b_n$ 都只是条件收敛, 而$\displaystyle \sum\limits_{n=2}^{\infty} c_n$ 发散的例子:
\[
a_n = b_n = \dfrac{(-1)^{n+1}}{\sqrt{n}}.
\]
\pause
我们接下来证明
\begin{prop}
若 $\displaystyle \sum\limits_{n=1}^{\infty} a_n, \sum\limits_{n=1}^{\infty} b_n, \sum\limits_{n=2}^{\infty} c_n$ 都收敛, 那么 $\displaystyle \sum\limits_{n=2}^{\infty} c_n = \left( \sum\limits_{n=1}^{\infty} a_n \right) \cdot \left( \sum\limits_{n=1}^{\infty} b_n \right).$
\end{prop}

\end{frame}

%------------------------------------------------

\begin{frame}
\frametitle{关于数项级数 Cauchy 乘积的扩展结论}

证明: 定义幂级数
\[
f(x) = \sum\limits_{n=1}^{\infty} a_n x^n, ~~
g(x) = \sum\limits_{n=1}^{\infty} b_n x^n, ~~
h(x) = \sum\limits_{n=1}^{\infty} c_{n+1} x^n.
\]
由 Abel 第二定理, $f,g,h$ 在 $[0, 1]$ 上一致收敛, 从而连续. 又由 Cauchy-Hadamard 定理 (或者 Abel 第一定理), $\forall ~ |x| < 1,$ $f(x),g(x),h(x)$ 绝对收敛, 于是有等式

\end{frame}

%------------------------------------------------

\end{document}
